\begin{proposition}\label{prop:2.2.8}
(2.14) 或 (2.15) 中定义的集合 $\pi^R_1((X,\sigma);x)$ 是
$\Aut^{\pm}(\tilde X)^{\mathrm{op}}$ 的子群：若 $\sigma_f=\Id_X$，则
$f\in\pi^R_1((X,\sigma);x)$ 全纯地作用在 $\tilde X$ 上；若 $\sigma_f=\sigma$，则它反
全纯地作用。存在良定义的群同态
\begin{equation}\label{eq:2.16}
\alpha:\ \pi^R_1((X,\sigma);x)\longrightarrow\{\Id_X;\sigma\}^{\mathrm{op}}\simeq\Z/2\Z,
\qquad f\longmapsto\sigma_f,
\end{equation}
其核为 $\ker\alpha=\{f:\tilde X\to\tilde X\mid p\circ f=p\}^{\mathrm{op}}=\pi_1(X;x)$。
特别地，$\pi_1(X;x)$ 是 $\pi^R_1((X,\sigma);x)$ 的正规子群。
\end{proposition}

我们将在后面看到 $\alpha$ 是满射，因而 $\pi_1(X;x)$ 在 $\pi^R_1((X,\sigma);x)$ 中的指数
为 $2$；见命题 2.2.12。

\begin{proposition}\label{prop:2.2.9}
设 $(X,\sigma)$ 是克莱因曲面，$q:(Y,\tau)\to(X,\sigma)$ 是实覆盖。回忆：给定
$y\in q^{-1}(\{x\})$，存在唯一映射 $t_y:\tilde X\to Y$ 使得 $t_y(\tilde x)=y$ 且
$q\circ t_y=p$，其中 $p:(\tilde X,\tilde x)\to(X,x)$ 是 $X$ 的一个固定万有覆盖。
(2.14) 中定义的群 $\pi^R_1((X,\sigma);x)$ 通过
\[
f\curvearrowright y:=(\tau_f^{-1}\circ t_y\circ f)(\tilde x)
\]
左作用于纤维 $q^{-1}(\{x\})$，其中当 $\sigma_f=\Id_X$ 时 $\tau_f=\Id_Y$，当
$\sigma_f=\sigma$ 时 $\tau_f=\tau$。此外，若 $\phi:(Y,\tau)\to(Y',\tau')$ 是
$(X,\sigma)$ 的两个实覆盖（设为 $q:(Y,\tau)\to(X,\sigma)$ 与
$q':(Y',\tau')\to(X,\sigma)$）之间的同态，则映射
$\phi_x:q^{-1}(\{x\})\to(q')^{-1}(\{x\})$ 是 $\pi^R_1((X,\sigma);x)$-等变的。
\end{proposition}

命题 2.2.9 实际上说明：$\pi^R_1((X,\sigma);x)$ 是定义在 $(X,\sigma)$ 的实覆盖范畴上的
「$x$ 处纤维」函子的自同构群。不过我们不会使用这一形式体系。

\begin{proof}[命题 2.2.8 的证明]
回忆 $\Aut^{\pm}(\tilde X)^{\mathrm{op}}$ 上的群结构为 $f\cdot g=g\circ f$。首先注意：
给定 $f\in\Aut^{\pm}(\tilde X)^{\mathrm{op}}$，满足 $p\circ f=\mu\circ p$ 的元素
$\mu\in\{\Id_X,\sigma\}$ 若存在则唯一。特别地，(2.16) 中的映射 $\alpha$ 是良定义的；
又因 $p$ 是全纯覆盖映射，若 $\sigma_f=\Id_X$ 则 $f$ 全纯，若 $\sigma_f=\sigma$ 则 $f$
反全纯。其次，若 $f,g\in\pi^R_1((X,\sigma);x)$，则
\[
p\circ(f\cdot g)=p\circ(g\circ f)=(p\circ g)\circ f
=(\sigma_g\circ p)\circ f=\sigma_g\circ(p\circ f)
=\sigma_g\circ(\sigma_f\circ p)=(\sigma_g\circ\sigma_f)\circ p
=(\sigma_f\cdot\sigma_g)\circ p,
\]
故 $f\cdot g\in\pi^R_1((X,\sigma);x)$，且由唯一性 $\sigma_{f\cdot g}=\sigma_f\cdot\sigma_g$。
同样，若 $f\in\pi^R_1((X,\sigma);x)$，则一方面 $p\circ(f\circ f^{-1})=p\circ\Id_{\tilde X}
=\Id_X\circ p=p$，另一方面 $(p\circ f)\circ f^{-1}=(\sigma_f\circ p)\circ f^{-1}$，
故 $p\circ f^{-1}=\sigma_f^{-1}\circ p$，这说明 $f^{-1}\in\pi^R_1((X,\sigma);x)$。因此
$\pi^R_1((X,\sigma);x)$ 是 $\Aut^{\pm}(\tilde X)^{\mathrm{op}}$ 的子群，且 $\alpha$ 是
群同态。由于
\[
\pi_1(X;x)=\{f:\tilde X\to\tilde X\mid p\circ f=\Id_X\circ p\}^{\mathrm{op}},
\]
$\pi_1(X;x)$ 正是 $\alpha$ 的核。
\end{proof}

\begin{proof}[命题 2.2.9 的证明]
下面的证明应与命题 2.2.5 的证明对照阅读。我们已在 (2.11) 与 (2.13) 中观察到，对一切
$f\in\pi^R_1((X,\sigma);x)$ 有
$f\curvearrowright y:=(\tau_f^{-1}\circ t_y\circ f)(\tilde x)\in q^{-1}(\{x\})$。此外有
$t_{f\curvearrowright y}=\tau_f^{-1}\circ t_y\circ f$。事实上，一方面按
$f\curvearrowright y$ 的定义 $(\tau_f^{-1}\circ t_y\circ f)(\tilde x)=f\curvearrowright y$；
另一方面
\[
\begin{aligned}
q\circ(\tau_f^{-1}\circ t_y\circ f)
&=(q\circ\tau_f^{-1})\circ(t_y\circ f)
=(\sigma_f^{-1}\circ q)\circ(t_y\circ f)\\
&=\sigma_f^{-1}\circ(q\circ t_y)\circ f
=\sigma_f^{-1}\circ p\circ f
=\sigma_f^{-1}\circ\sigma_f\circ p=p,
\end{aligned}
\]
故由映射 $t_{f\curvearrowright y}$ 的唯一性得
$\tau_f^{-1}\circ t_y\circ f=t_{f\curvearrowright y}$。\begin{annote}原文此处印作
$\tau_f\circ t_y\circ f=t_{f\curvearrowright y}$，漏掉了逆号；按上文计算与下文的用法，
应为 $\tau_f^{-1}\circ t_y\circ f=t_{f\curvearrowright y}$。译文按后者写出（参见疑点清单
Q006）。\end{annote}特别地，这蕴涵 $\pi^R_1((X,\sigma);x)$ 左作用于 $q^{-1}(\{x\})$。
事实上，注意到 $\sigma_{f\cdot g}=\sigma_f\cdot\sigma_g$ 按 $\tau_f$ 的定义蕴涵
$\tau_{f\cdot g}=\tau_f\cdot\tau_g$，于是
\[
\begin{aligned}
(f\cdot g)\curvearrowright y
&=\bigl(\tau_{f\cdot g}^{-1}\circ t_y\circ(f\cdot g)\bigr)(\tilde x)\\
&=\bigl((\tau_f\cdot\tau_g)^{-1}\circ t_y\circ(f\cdot g)\bigr)(\tilde x)\\
&=\bigl((\tau_g\circ\tau_f)^{-1}\circ t_y\circ(g\circ f)\bigr)(\tilde x)\\
&=\bigl(\tau_f^{-1}\circ(\tau_g^{-1}\circ t_y\circ g)\circ f\bigr)(\tilde x)
=\bigl(\tau_f^{-1}\circ t_{g\curvearrowright y}\circ f\bigr)(\tilde x)\\
&=f\curvearrowright(g\curvearrowright y).
\end{aligned}
\]
同样，$\phi_x:q^{-1}(\{x\})\to(q')^{-1}(\{x\})$ 的等变性来自图
\[
\begin{tikzcd}[column sep=large, row sep=large]
 & \tilde X \arrow[dl, "t_y"'] \arrow[dr, "t_{\phi(y)}"] & \\
 (Y,\tau) \arrow[dr, "q"'] \arrow[rr, "\phi"] & & (Y',\tau') \arrow[dl, "q'"]\\
 & (X,\sigma) &
\end{tikzcd}
\]
的可交换性。事实上，我们在 (2.12) 中已证 $\phi\circ t_y=t_{\phi(y)}$。于是对一切
$f\in\pi^R_1((X,\sigma);x)$，
\[
\begin{aligned}
\phi(f\curvearrowright y)
&=\phi\bigl((\tau_f^{-1}\circ t_y\circ f)(\tilde x)\bigr)
=\bigl((\phi\circ\tau_f^{-1})\circ t_y\circ f\bigr)(\tilde x)\\
&=\bigl((\tau'_f)^{-1}\circ\phi\circ t_y\circ f\bigr)(\tilde x)\\
&=\bigl((\tau'_f)^{-1}\circ(\phi\circ t_y)\circ f\bigr)(\tilde x)
=\bigl((\tau'_f)^{-1}\circ t_{\phi(y)}\circ f\bigr)(\tilde x)\\
&=f\curvearrowright\phi(y),
\end{aligned}
\]
故 $\phi_x:=\phi|_{q^{-1}(\{x\})}$ 确实是 $\pi^R_1((X,\sigma);x)$-等变的。
\end{proof}

现在我们可以给群 $\pi^R_1((X,\sigma);x)$ 一个名称了。$\pi^R_1((X,\sigma);x)$ 的另一种
记号或许可以是 $(\Aut^\sigma(\tilde X/X))^{\mathrm{op}}$，但正如复曲线的基本群的情形，
它没有把基点考虑在内，而这在以后会造成问题（见命题 2.2.13 与命题 2.2.17）。在给出
本讲义核心概念的定义之前，最后再作一点说明：关于克莱因曲面的基本群，还有若干更具
概念性的理解方式（例如进一步展开范畴的观点、谈论纤维函子的自同构，或者考虑轨形并
使用轨形基本群的一般理论）；这里我们采取相当「动手算」的做法，只考虑克莱因曲面的实
覆盖与底黎曼面的拓扑万有覆盖。我们希望这能作为对该主题的平易入门，并让读者渴望了解
这个丰富课题的更多内容。

\begin{definition}[实曲线的基本群]\label{def:2.2.10}
设 $(X,\sigma)$ 是克莱因曲面，$x\in X$。固定万有覆盖
$p:(\tilde X,\tilde x)\to(X,x)$。克莱因曲面 $(X,\sigma)$ 的\emph{基本群}是群
\[
\pi^R_1((X,\sigma);x):=
\bigl\{f:\tilde X\to\tilde X\ \bigm|\ \exists\sigma_f\in\{\Id_X;\sigma\},\
p\circ f=f\circ\sigma_f\bigr\}^{\mathrm{op}},
\]
其群结构为 $(f\cdot g):=g\circ f$。
\end{definition}

正如前面简略提到的，更具概念性的定义是直接把 $\pi^R_1((X,\sigma);x)$ 取为
$(X,\sigma)$ 的实覆盖范畴上「$x$ 处纤维」函子的自同构群；但由于万有覆盖
$p:(\tilde X,\tilde x)\to(X,x)$ 的存在，这最终会与刚才给出的定义一致。应当注意，
$X$ 的万有覆盖 $\tilde X$ 一般并不是实覆盖：$X$ 在黎曼面范畴中有连通万有覆盖，但
$(X,\sigma)$ 在克莱因曲面范畴中一般没有连通万有覆盖（看出这一点的一种办法是使用黎曼
面的单值化定理，并例如证明：Poincaré 上半平面上的实结构必被 $\mathrm{PSL}(2;\R)$ 中
某元素共轭于 $z\mapsto-\bar z$，因而特别地其实点集连通；然而存在亏格 $g\ge2$ 的克莱因
曲面，其上带有无实点的实结构，从而在克莱因曲面范畴中不能被上半平面覆盖）。这也意味着
一般不能指望 $\pi^R_1((X,\sigma);x)$ 通过实覆盖的自同构作用在 $\tilde X$ 上。不过，
$\pi^R_1((X,\sigma);x)$ 通过或全纯或反全纯的变换作用在 $\tilde X$ 上；换言之，
$\pi^R_1((X,\sigma);x)$ 是 $(\Aut^{\pm}(\tilde X))^{\mathrm{op}}$ 的子群。

下面这个引理将在本小节余下部分多次用到。
